% Módulo 02d — Trabalhos relacionados (resumos como fonte, R773+)
Com teoria e ferramenta postas, situamos o trabalho na literatura --- declarando o que foi apenas lido.
\section{Trabalhos relacionados}
\textit{Calibração por temperatura.} A literatura registra a temperatura como parâmetro de nitidez e calibração: relações entre focal loss, temperature scaling e Scorings próprios, com o hiato de generalização gerando sobreconfiança no teste \cite{komisarenko2024}; aspectos práticos e teóricos em classificação e robustez adversarial \cite{xuan2025}; e escalonamento adaptativo para modelos de linguagem \cite{xie2024}. Nossa T10 contribui a forma elementar verificada por máquina, sem limites assintóticos. \textit{Ponto flutuante.} Aproximação correta de IEEE-754 para verificação de programas, com arredondamentos, infinitos, NaNs e zeros sinalizados \cite{bagnara2019}; certificação de robustez via Lipschitz sob execução float, com lacunas float-relevantes para garantias \cite{murray2026}. Contribuímos T11 e teste empírico de propriedades, sem reivindicar certificação total. \textit{Benchmarks de agentes.} Avaliação multidimensional de LLMs como agentes em ambientes interativos \cite{agentbench2023}; lacuna de perspectiva em prompting de orquestração multiagente \cite{perspectivegap2026}; modelagem de recompensa para orquestradores \cite{orchrm2026}; arcabouço com coordenação auto-emergente inspirada em atenção \cite{orchestrator2025} --- convergência independente digna de nota, não validação. Nenhum foi executado aqui. (Síntese limitada aos resumos; leitura integral futura.)
\section{Fichamento por obra (o que cada uma informa)}
\textit{Calibração:} relações entre perda focal, scaling e propriedade \cite{komisarenko2024} informam T10; prática e teoria em classificação e robustez \cite{xuan2025} informam a leitura da Tabela~\ref{tab:temp}; escalonamento adaptativo para linguagem \cite{xie2024} informa limites do que não fizemos. \textit{Floats:} aproximação correta de IEEE-754 \cite{bagnara2019} informa o desenho do fuzz; certificação via Lipschitz sob floats \cite{murray2026} informa o que faltaria para certificar. \textit{Benchmarks:} avaliação multidimensional \cite{agentbench2023}, lacuna de perspectiva \cite{perspectivegap2026}, recompensa de orquestração \cite{orchrm2026} e coordenação atencional \cite{orchestrator2025} informam o protocolo externo pendente. \textit{Verificação:} base axiomática \cite{hoare1969}, Lean \cite{moura2021}, Mathlib \cite{mathlib2020}, Z3 \cite{demoura2008} e SymPy \cite{meurer2017} informam o núcleo. \textit{Método:} experimentação \cite{basili1986,wohlin2012}, casos \cite{runeson2009}, relato \cite{page2021}, valores-\textit{p} \cite{wasserstein2016}, estimação \cite{cumming2014}, bootstrap \cite{efron1979}, reprodutibilidade \cite{osc2015,collberg2016,ioannidis2005}, dados e modelos \cite{gebru2021,mitchell2019}, FAIR \cite{wilkinson2016} informam desenho, relato e trilha.
\par Leitura por obra: calibração informa T10; ponto flutuante informa T11 e o fuzz; benchmarks informam o protocolo externo pendente; verificação informa o núcleo Lean; método informa desenho, relato e estatística; reprodutibilidade informa a trilha; dados e modelos documentados informam manifestos. Nenhuma obra foi executada: são referências de contexto, não resultados reivindicados.
\par Comentário teórico: cada família acima é relacionada na Tabela~\ref{tab:qkv} e no texto a uma lacuna do trabalho (temperatura a T10, ponto flutuante a T11 e ao fuzz, benchmarks ao protocolo externo pendente); registra-se que esses trabalhos foram apenas lidos em resumo, sem execução, e não integram os resultados reivindicados.
\section{Leitura por obra (resumos como fonte)}
\textit{Atenção e linguagem.} \cite{vaswani2017} informa a forma herdada; \cite{luong2015}, que a forma admite variantes; \cite{devlin2019}, o alcance da arquitetura, não usado aqui; \cite{dai2019}, limites do que não foi feito; \cite{wolf2020}, o ferramental consolidado do campo. \textit{Benchmarks e avaliação.} \cite{wang2018}, \cite{rajpurkar2016} e \cite{deng2009} fixaram o padrão de avaliação por tarefas aqui adaptado em escala sintética; \cite{agentbench2023}, \cite{perspectivegap2026}, \cite{orchrm2026} e \cite{orchestrator2025} definem o palco externo ainda não pisado. \textit{Calibração e ponto flutuante.} \cite{komisarenko2024}, \cite{xuan2025} e \cite{xie2024} informam T10 e sua fronteira; \cite{bagnara2019} e \cite{murray2026} informam T11, o fuzz e o que faltaria certificar; fundamentos em \cite{muller2018,goldberg1991,harris2020,virtanen2020,meurer2017,demoura2008}. \textit{Método e auditoria.} Desenho e relato em \cite{basili1986,wohlin2012,runeson2009,page2021}; relato por intervalos em \cite{wasserstein2016,cumming2014,efron1979,gneiting2007}; trilha em \cite{osc2015,collberg2016,ioannidis2005,gebru2021,mitchell2019,wilkinson2016}; auditoria em \cite{fagan1976,mccabe1976,chidamber1994,parnas1986}; linhagem da orquestração em \cite{smith1980,erman1980}; verificação dedutiva em \cite{hoare1969,moura2021,mathlib2020}. Leitura limitada aos resumos; nenhuma obra foi executada.
\section{Síntese: matriz de cobertura deste trabalho}
Numa matriz formal $\times$ empírico $\times$ auditável: formal coberto para o núcleo (11/11 verificados); empírico coberto internamente (bancada, ablações, fuzz, correlações); auditável coberto (axiomas, alegações, referências, pinos, sincronia); externo descoberto (sem head-to-head, sem reprodução independente, sem Q/K/V formal completo). A novidade reivindicável é composicional --- verificação, avaliação e auditoria amarradas sobre um mesmo objeto com trilha local de reprodução --- não a de cada peça isolada.
